\section{\texttt{vm\_area\_t} 结构体}

VMA 系统中最基本的数据结构是 \texttt{vm\_area\_t}——描述一段虚拟内存区域：

\codefile{src/include/kernel/vma.h}
\begin{lstlisting}[style=cstyle]
typedef struct vm_area {
    uint32_t    start;      /* 起始虚拟地址（包含，页对齐） */
    uint32_t    end;        /* 结束虚拟地址（不包含，页对齐） */
    uint32_t    flags;      /* VMA_READ | VMA_WRITE | VMA_EXEC */
    const char* name;       /* 人类可读名称, 如 "direct-map", "kheap" */
} vm_area_t;
\end{lstlisting}

约定与 Linux 一致：\texttt{[start, end)} 左闭右开，\texttt{start} 和 \texttt{end} 都必须页对齐，
不同 VMA 之间不允许地址重叠。

\begin{figure}[H]
  \centering
  % figures/ch11/fig01-vm-area-layout.tex — auto-extracted from chapter source
\begin{tikzpicture}[
  cell/.style={draw, minimum height=0.7cm, font=\ttfamily\small, anchor=west},
]
  \node[font=\scriptsize, anchor=south, text=blue!60!black] at (0.2, 0.9)
    {\texttt{vm\_area\_t*}};
  \draw[-{Stealth[length=4pt]}, thick, blue!60!black] (0.2, 0.9) -- (0.2, 0.7);
  \node[cell, minimum width=2.4cm, fill=orange!12]  at (0, 0)   {\small start};
  \node[cell, minimum width=2.4cm, fill=orange!12]  at (2.4, 0) {\small end};
  \node[cell, minimum width=2.4cm, fill=yellow!15]  at (4.8, 0) {\small flags};
  \node[cell, minimum width=2.4cm, fill=green!12]   at (7.2, 0) {\small name};
  \draw[{|}-{|}, thick] (0, -0.6) -- (2.4, -0.6);
  \node[font=\scriptsize, below] at (1.2, -0.6) {4\,B};
  \draw[{|}-{|}, thick] (2.4, -0.6) -- (4.8, -0.6);
  \node[font=\scriptsize, below] at (3.6, -0.6) {4\,B};
  \draw[{|}-{|}, thick] (4.8, -0.6) -- (7.2, -0.6);
  \node[font=\scriptsize, below] at (6.0, -0.6) {4\,B};
  \draw[{|}-{|}, thick] (7.2, -0.6) -- (9.6, -0.6);
  \node[font=\scriptsize, below] at (8.4, -0.6) {4\,B};
  \draw[{|}-{|}, thick] (0, 0.9) -- (9.6, 0.9);
  \node[font=\scriptsize, above] at (4.8, 0.9) {sizeof = 16\,B};
\end{tikzpicture}

  \caption{\texttt{vm\_area\_t} 内存布局（32 位平台）}
  \label{fig:vm-area-layout}
\end{figure}

注意：\texttt{vm\_area\_t} 是\emph{公开的只读视图}，不包含后端内部管理字段（如红黑树颜色、
父/子指针等）。这保持了接口与实现的解耦——调用者只读 \texttt{vm\_area\_t}，
后端的内部节点可以包含任意附加字段。

% ============================================================================

